\documentclass[instructions.tex]{subfiles}
\begin{document}
 
\chapter{Nos devoirs envers Jésus-Christ dans la sainte Eucharistie}
 
Jésus-Christ est réellement dans ce mystère; Jésus-Christ s’y sacrifie; Jésus-Christ s’y donne.

\primo Jésus-Christ est dans ce mystère. Nous devons le croire; nous devons l’y adorer; nous devons le visiter. 1. A qui croirons-nous, si nous ne croyons pas à un Dieu qui parle? N’est-il donc pas assez puissant pour faire ce qu’il dit? Ceci est mon corps, qui sera livré pour vous, dit Jésus-Christ; ce calice est mon sang qui sera répandu pour vous. Or, ce n’est pas la figure de son corps qui a été livré pour nous, mais son véritable corps; ce n’est pas la figure de son sang, mais son vrai sang qui a été répandu: son vrai corps et son sang sont donc dans ce mystère. Je n’ai pas besoin du témoignage de mes yeux pour m’assurer de la présence de Jésus-Christ sur l’autel, disait saint Louis; la foi me suffit, je le crois avec plus de fermeté que si je le voyais en personne.

2. Puisque Jésus-Christ est présent dans ce mystère, quels sentimens de respect et d’adoration ne lui devons-nous pas! Si nous l’avions vu lorsqu’il était sur la terre, avec quelle profonde humilité aurions-nous approché de sa personne adorable! Lui en devons-nous moins, parce que nous ne le voyons pas? Les démons tremblent en sa présence; des millions d’esprits célestes remplissent ses temples, et l’adorent avec une sainte frayeur: et nous, misérables! nous osons y paraître avec dissipation! On tremble devant un roi mortel, et nous sommes sans respect devant un Dieu!  Ne savons-nous pas que, s’il est dans ce mystère un Dieu caché, il n’en est pas moins un Dieu terrible?

3. La reconnoissance, l’amour et la confiance doivent nous engager à lui rendre de fréquentes visites. C’est pour notre amour qu’il habite parmi nous: il nous invite à aller à lui; pourquoi nous éloignons-nous d’un Dieu qui nous appelle? Lorsqu’il était sur la terre, les peuples en foule accouraient à lui; s’il allait dans le désert, les peuples le suivaient; s’il entrait dans les villes, s’il allait au temple, les peuples l’accompagnoient. C’est le même Jésus-Christ qui habite parmi nous, et nous le laissons seul. Nous n’avons point d’ami plus tendre; il est notre Dieu, notre Sauveur et notre Père, et nous vivons avec lui comme s’il nous était étranger! Quelle ingratitude!

Apprenons des courtisans qui font sans cesse la cour à leur prince, ce que nous devons à Jésus-Christ. Devons-nous faire moins pour le roi du ciel, que tant de flatteurs mercenaires ne font pour les rois de la terre?  N’est-il pas honteux qu’on dise de nous ce que saint Jean-Baptiste reprochait aux Juifs? Jésus-Christ est au milieu de vous, et vous ne le connoissez pas, vous l’abandonnez\footnote{Medius autem vestrûm stetit quem vos nescitis. Joan. 1.}.

Après tout, c’est lui qui doit nous juger: sans lui nous ne pouvons être sauvés: ainsi, l’état de notre âme, nos tentations, notre misère, notre salut, doivent nous faire sentir le besoin continuel que nous avons de son secours et de sa miséricorde. Si vous êtes assidu à le visiter dans le Saint-Sacrement, soyez assuré que vous ne sortirez jamais de sa présence sans recevoir quelques nouvelles grâces. Le démon, la paresse, les compagnies, le jeu, vous éloigneront d’une si sainte pratique; les écouterez-vous au préjudice de ce que vous devez à votre Sauveur et à vous-même? Si vos occupations ne vous permettent pas de le visiter dans son temple, du moins pensez souvent à lui, en l’ado- rant, en implorant son secours.

\secundo Jésus-Christ se sacrifie dans ce mystère. Le Fils de Dieu, dit saint Paul, s’est abaissé jusqu’au néant, en se faisant homme; mais il s’est humilié bien plus profondément, en se faisant victime sur la croix. Or, tous les jours il continue de se faire victime, en s’offrant en sacrifice à son Père. Considérez, lorsque vous assistez à la messe, pourquoi Jésus-Christ s’y offre en sacrifice. C’est pour faire ce que nous devons faire nous-mêmes, et ce que tous les hommes ne peuvent faire dignement sans lui; c’est pour y adorer son Père, et le remercier pour nous; c’est pour lui demander pour nous des grâces, et apaiser la justice de son Père sur nous.

Oh! combien est agréable à Dieu cette victime humiliée! Dieu pourrait-il souffrir tant de crimes sur la terre, s’il ne voyait pas sur nos autels son Fils qui se sacrifie pour les expier? Assistons à ce sacrifice adorable, en nous unissant à cette victime sainte pour rendre nos hommages au Tout-Puissant, pour obtenir le pardon et sa miséricorde; assistons-y souvent, puisqu’une seule messe rend plus de gloire à Dieu que les bonnes œuvres de tous les justes; assistons-y avec les mêmes sentimens de douleur et de reconnoissance que nous aurions eus, si nous eussions assisté à sa mort sur le Calvaire. Mourons en esprit avec lui, en faisant mourir nos passions\footnote{Eamus et nos, ut moriamur cum eo. Joan. 11. 13.}.

\tertio Jésus-Christ se donne à vous dans ce mystère.  Quelle bonté dans un Dieu! La mère la plus tendre a-t-elle jamais fait pour un fils ce que Jésus-Christ fait ici pour nous? Il veut être la nourriture de notre âme pendant la vie; et, pour marquer encore plus sa charité, il veut lui-même nous servir de viatique et de sauve-garde à la mort.

Oh! que le voyage de cette vie pour aller au ciel est bien plus difficile que celui du prophète Élie à la montagne d’Horeb! Tout y est rempli d’écueils et d’ennemis. Dieu envoya à ce prophète un pain matériel pour le fortifier; mais Jésus-Christ nous donne un pain céleste, qui est son propre corps. Prenons ce pain divin avec foi; nous serons terribles au démon même, dit saint Chrysostome. C’est surtout à la mort que nous ressentirons avec consolation ses effets heureux dans le sacré viatique. Le démon redoublera alors ses efforts pour nous perdre; mais qu’aurons-nous à craindre, lorsque par la communion nous aurons Jésus-Christ avec nous?

Concluons avec le Prophète, qu’il n’y a point de nation qui ait des dieux comme le nôtre, qui s’abaisse ainsi jusqu’à nous\footnote{Deut. 4. 7.}. Les paiens qu’on instruisait dans nos mystères avaient bien raison de s’écrier: Oh! que le Dieu des chrétiens est bon! qu’il mérite d’être aimé!  Nous sommes bien malheureux, dit un saint évêque, si, ayant autour de nous les charbons ardens de la charité d’un Dieu, nos cœurs, par un prestige du démon, sont tout de glace pour ce Dieu d’amour\footnote{Homo, tot congestis carbonibus, miraculo diabolico friget ad Deum. Guill. Paris.}. Nous sommes étrangement ennemis de nous-mêmes, si nous nous éloignons de lui.
 
\section{Résolutions}
 
\primo Je ne paraîtrai plus devant Jésus-Christ qu’avec un profond respect.

\secundo Je le visiterai tous les jours, et j’assisterai de même à la sainte messe, autant que mes occupations et ma santé me le permettront.

\tertio Et souvent dans la journée, je porterai mon cœur aux pieds des autels, surtout quand je me sentirai pressé par quelque tentation.

\prayer{O Jésus! je vous adore dans l’auguste sacrement de votre amour. Ah! pardonnez-moi mes ingratitudes passées.}
 
\end{document}
